\documentclass[11pt]{article}
\usepackage[margin=0.9in]{geometry}
\usepackage{amsmath,booktabs,hyperref}
\title{Flow-seed visualization and a local infeasibility bound}
\author{Pan Dynamics Collision Avoidance Research}
\date{October 2, 2026}
\begin{document}
\maketitle
The exact fresh-flow initialization at 0.85 s in the 15 m/s turning-crossing
case already overlaps the target at absolute time 1.65 s. A necessary linear
constraint bound identifies why the original local problem cannot repair it.
This is an offline visualization and diagnosis, not a controller modification.

\section{What is visualized}
The saved fresh-flow anchor is a nonlinear vehicle rollout of the flow-guided
feedback inputs, including seed input shaping and the completion behavior.
It is not an ideal geometric streamline, the optimized affine trajectory,
or the subsequently executed closed-loop trajectory. Both vehicles are shown
at matching absolute times. The given path is the horizontal line; there are
no road-boundary constraints. Current time is 0.85 s, the horizon ends at
4.65 s, and the 77 stored nodes are spaced by 0.05 s. The primary portion ends
at 1.65 s, where the hard completion portion starts.

The ego initially has position $(11.9112,-0.0204)$ m and longitudinal speed
12.8219 m/s. The target initially has position $(24.8104,7.1577)$ m and speed
9.25 m/s. At 1.65 s the seed ego position is $(21.0017,1.7466)$ m, while the
target position is $(24.0039,-0.5012)$ m. Their rectangles overlap with signed
SAT gap $-0.290161$ m. This is a sampled geometric overlap measure, not crash
penetration deformation. The visualization includes a current-state overview,
a same-time overlap close-up, and sampled gap history. The local HTML viewer
supports time scrubbing, playback, a jump to overlap, and full-horizon view.

\section{Why the original local problem is infeasible}
The original correction box has input increments bounded by
$|\Delta\delta|\le0.075$ rad and $|\Delta b|\le0.125$, together with the
configured state boxes. These are numerical local bounds. Reconstructing them
from the saved enlarged problem leaves the same anchor and affine equations.
For one corner-separation row at stage 17, the buffer-inclusive margin is
\[
 g(d)=-0.390161412+j^Td,\qquad g(d)\ge0.
\]
Stage 17 is hard, so no safety slack can relax this row. To bound what the
local model can accomplish, solve
\[
 \overline g=\max_d g(d)\quad\text{subject to}\quad
 Ed=r,\quad \ell\le d\le u.
\]
Here $Ed=r$ contains the affine dynamics, and $[\ell,u]$ contains the saved
variable bounds reconstructed at input scale one. All other inequality rows
and the terminal cone are omitted. This is a relaxation of the actual problem,
so its maximum bounds the margin attainable in the complete local problem.
MATLAB \texttt{linprog} returns
\[
 \overline g=-0.048075525\ \mathrm{m}<0.
\]
The row needs 0.390161 m of improvement, but at most 0.342086 m is available
in this relaxed calculation. Thus this one hard row is already impossible
inside the original local problem. The numerical optimizer returns flag 1,
with affine-equality residual $1.31\times10^{-9}$ and zero box violation.
The row coefficients and bound match row 225 of the saved formulation.
A separate solve of all linear inequalities without the terminal cone also
returns infeasible. The endpoint cone is therefore not needed to explain
infeasibility of this particular fresh-flow problem.

This is a numerical necessary-condition diagnosis of the local affine problem,
not a global physical infeasibility certificate. The 4.8 cm deficit is relative
to the configured 10 cm margin, not a claim that every physically possible
nearby maneuver must collide. With zero requested buffer this row alone would
no longer establish infeasibility; the other rows have not been certified
feasible under such a change. No margin was changed in this task.

\section{Implication and reproducibility}
An unsafe seed is allowed by design. The specific failure is that this seed
and the available local correction set cannot jointly satisfy a hard future
separation row. A useful initialization need not be safe, but the solver needs
a way to reach a safe trajectory from it. The prior diagnosis showed that a
larger correction box finds a near-zero-slack affine result with inaccurate
nonlinear motion; simply enlarging the box is not established as a remedy.

The companion directory \path{report/FLOW_SEED_VISUALIZATION_20261002}
contains the exact seed export, row bound, source hashes, and reproduction
scripts. Run \texttt{exportFlow.m} against its declared original MATLAB
snapshots; run \texttt{render.py DATA\_DIRECTORY OUTPUT\_DIRECTORY} to render
outside the repository. All 77 independently computed rectangle distances
match MATLAB within $10^{-10}$ m. The HTML was rendered in Chromium and its
screenshot visually reviewed, then opened in the desktop browser. No complete
scenario campaign or controller unit suite was rerun. Generated graphics,
HTML and PDF are retained in the external diagnostic cache and exported to
the research archive. Production code and unrelated working-tree changes
remain untouched.
\end{document}
